\chapter{Engineering Standards and Design Challenges}

This chapter details compliance with international software, security, and metadata engineering standards, evaluates societal and environmental impacts, provides a comparative financial budget analysis, and presents the BAETE Complex Engineering Problem (EP1--EP7), Knowledge Profile (K1--K8), and Complex Engineering Activity (EA1--EA5) accreditation mappings.

\section{Compliance with the Standards}

\subsection{Software Standards}
\begin{itemize}
    \item \textbf{ISO/IEC 25010 (Software Product Quality)}: Enforced to assess functional suitability, reliability under concurrent booking requests, performance efficiency, and usability across desktop and mobile form factors.
    \item \textbf{REST Architecture (RFC 7231 / RFC 7235)}: Stateless client-server communication using standardized HTTP verbs (\texttt{GET}, \texttt{POST}, \texttt{PATCH}, \texttt{DELETE}) and uniform JSON resource representations.
    \item \textbf{JSON Specification (RFC 8259)}: Universal serialization for all API request payloads and responses, ensuring cross-platform data interchange.
    \item \textbf{W3C WCAG 2.1 Level AA}: High-contrast interface components, responsive touch targets, and semantic HTML accessibility for diverse student populations.
\end{itemize}

\subsection{Security and Communication Protocols}
\begin{itemize}
    \item \textbf{TLS 1.3 (RFC 8446)}: Enforcing end-to-end cryptographic encryption for all student sessions, preventing man-in-the-middle credential interception.
    \item \textbf{JSON Web Tokens (RFC 7519)}: Compact, URL-safe cryptographic claims with expiration headers stored exclusively in secure \texttt{HttpOnly} cookies to mitigate Cross-Site Scripting (XSS).
    \item \textbf{QR Code Standard (ISO/IEC 18004:2015)}: Standardized error-correcting 2D barcodes for rapid seat arrival verification via camera scanning.
    \item \textbf{Dublin Core Metadata Standard (ISO 15836)}: Standardized bibliographic schemas (Title, Author, Publisher, Year, ISBN, Subject Classification) for universal catalog interoperability.
\end{itemize}

\section{Impact on Society, Environment and Sustainability}

\subsection{Impact on Life and Student Well-Being}
By removing manual seat contention, chaotic desk reservation habits, and extensive physical book search latencies, the system significantly lowers student stress levels during high-pressure examination cycles. Fair slot allocations guarantee equitable access to university resources for all enrolled students.

\subsection{Impact on Society and Environmental Sustainability}
The platform replaces physical paper sign-in sheets, paper reservation cards, and printed checkout slips with digital verification passes. On a large campus with over 20,000 active students, this saves hundreds of reams of paper annually, actively contributing to university carbon footprint reduction and promoting green campus initiatives.

\subsection{Ethical Aspects and Data Privacy}
All user authentication records are protected using industry-standard Bcrypt hashing. Student reading history and seat usage records are secured behind strict Role-Based Access Control (RBAC), preventing unauthorized patron tracking or personal profiling.

\subsection{Sustainability and Maintenance Plan}
The application is built on open-source, non-proprietary frameworks (Node.js, Next.js, PostgreSQL, Linux) without expensive vendor lock-in. Containerized Docker deployment guarantees long-term maintainability, easy horizontal scaling across campus branches, and zero recurring commercial software licensing fees.

\section{Project Management and Financial Analysis}
Table~\ref{tab:budget_comparison} presents a comprehensive comparative financial analysis between the proposed in-house open-source architecture and proprietary commercial enterprise alternatives (such as Ex Libris Alma or proprietary room booking suites).

\begin{table}[H]
\centering
\caption{Proposed Budget vs. Commercial Alternate Budget Analysis}
\label{tab:budget_comparison}
\small
\begin{tabularx}{\textwidth}{>{\raggedright\arraybackslash}p{3.8cm} >{\centering\arraybackslash}p{2.6cm} >{\centering\arraybackslash}p{2.6cm} >{\raggedright\arraybackslash}X}
\toprule
\textbf{Cost Component / Item} & \textbf{Proposed Budget (In-House)} & \textbf{Commercial Alternate} & \textbf{Selection Rationale} \\
\midrule
Software Licensing Fees & 0 BDT (Open-Source) & 350,000 BDT / yr & Complete ownership with zero recurring per-seat user license fees. \\
Cloud Hosting \& Storage & 30,000 BDT / yr & 120,000 BDT / yr & Standard Linux VPS and managed PostgreSQL provide optimal price-to-performance. \\
Entry Check-In Tablets & 24,000 BDT & 95,000 BDT & Web-based responsive camera UI avoids costly proprietary turnstile scanners. \\
Barcode / QR Scanners & 8,500 BDT & 35,000 BDT & Universal USB/Bluetooth optical scanners compatible with standard web APIs. \\
Maintenance \& Upgrades & 15,000 BDT / yr & 80,000 BDT / yr & Maintained internally by university IT staff and student engineering teams. \\
\midrule
\textbf{Total 1st Year Cost} & \textbf{77,500 BDT} & \textbf{680,000 BDT} & \textbf{88.6\% Institutional Cost Savings} \\
\bottomrule
\end{tabularx}
\end{table}

\section{Complex Engineering Problem (CEP)}

\subsection{Complex Problem Solving Mapping}
Table~\ref{tab:ep_mapping} summarizes the project mapping to the BAETE Complex Engineering Problem (EP1--EP7) attributes.

\begin{table}[H]
\centering
\caption{Mapping with Complex Engineering Problem Attributes (BAETE Table 5.1)}
\label{tab:ep_mapping}
\small
\begin{tabularx}{\textwidth}{*{7}{>{\centering\arraybackslash}X}}
\toprule
\textbf{EP1} & \textbf{EP2} & \textbf{EP3} & \textbf{EP4} & \textbf{EP5} & \textbf{EP6} & \textbf{EP7} \\
\midrule
Depth of Knowledge & Conflicting Requirements & Depth of Analysis & Familiarity of Issues & Applicable Codes & Stakeholder Involvement & Interdependence \\
\midrule
\checkmark & \checkmark & \checkmark & \checkmark & \checkmark & \checkmark & \checkmark \\
\bottomrule
\end{tabularx}
\end{table}

\subsubsection{Rationale for EP1 (Depth of Knowledge)}
Resolving real-time concurrency, multi-tier spatial book location indexing, and cryptographic verification requires foundational and specialist knowledge spanning database concurrency control (ACID transactions), cryptographic hash functions, relational schema normalization, and full-stack software engineering.

\subsubsection{Rationale for EP2 (Range of Conflicting Requirements)}
The system balances conflicting technical constraints: sub-50\,ms search latency versus multi-table spatial joins, open student access versus tight administrative lending control, and rapid seat booking versus strict physical arrival verification windows.

\subsubsection{Rationale for EP3 (Depth of Analysis)}
Formalizing the mathematical probability of seat collision ($P(\text{Collision})=0$), dynamic slot time-window state machines, and relational constraint dependencies required deep algorithmic analysis.

\subsubsection{Rationale for EP4 (Familiarity of Issues)}
Addressing simultaneous registration surges, sudden phantom seat reservations, and unindexed physical library shelves represents non-standard, institution-specific engineering challenges.

\subsubsection{Rationale for EP5 (Extent of Applicable Codes)}
Strict compliance with ISO/IEC 25010, RFC 7231 (REST), RFC 7519 (JWT), and W3C WCAG accessibility standards.

\subsubsection{Rationale for EP6 (Extent of Stakeholder Involvement)}
Requires active coordination and continuous reconciliation between students, librarians, faculty members, and university IT administrators.

\subsubsection{Rationale for EP7 (Interdependence)}
High technical interdependence exists across the Next.js presentation layer, Express RESTful microservices, Prisma ORM layer, PostgreSQL database constraints, and camera QR scanning hardware.

\subsection{Mapping with Knowledge Profile}
Table~\ref{tab:knowledge_profile} details the BAETE Knowledge Profile (K1--K8) mappings.

\begin{table}[H]
\centering
\caption{Mapping with Knowledge Profile (BAETE Table 5.2)}
\label{tab:knowledge_profile}
\small
\begin{tabularx}{\textwidth}{*{8}{>{\centering\arraybackslash}X}}
\toprule
\textbf{K1} & \textbf{K2} & \textbf{K3} & \textbf{K4} & \textbf{K5} & \textbf{K6} & \textbf{K7} & \textbf{K8} \\
\midrule
Natural Science & Mathematics & Engineering Fundamentals & Specialist Knowledge & Engineering Design & Engineering Practice & Comprehension & Research Literature \\
\midrule
\checkmark & \checkmark & \checkmark & \checkmark & \checkmark & \checkmark & \checkmark & \checkmark \\
\bottomrule
\end{tabularx}
\end{table}

\begin{itemize}
    \item \textbf{K1 (Natural Science)}: Optical barcode/QR light reflection principles and camera sensor data capture.
    \item \textbf{K2 (Mathematics)}: Set theory, relational algebra, probability modeling of resource contention, and cryptographic hashing.
    \item \textbf{K3 (Engineering Fundamentals)}: Computer organization, operating system processes, networking, and relational database management.
    \item \textbf{K4 (Specialist Knowledge)}: Advanced web architecture, asynchronous event loops, ORM query optimization, and RBAC security.
    \item \textbf{K5 (Engineering Design)}: Systematic design from stakeholder requirements to DFD, ERD, and full-stack software deployment.
    \item \textbf{K6 (Engineering Practice)}: Industry-standard toolchains: TypeScript, Next.js, Prisma, PostgreSQL, Git, and Docker.
    \item \textbf{K7 (Comprehension)}: Appreciation of societal convenience, academic equity, and paperless environmental sustainability.
    \item \textbf{K8 (Research Literature)}: Systematic review and synthesis of peer-reviewed smart library research papers.
\end{itemize}

\subsection{Mapping with Complex Engineering Activities}
Table~\ref{tab:ea_mapping} maps the project to Complex Engineering Activities (EA1--EA5).

\begin{table}[H]
\centering
\caption{Mapping with Complex Engineering Activities (BAETE Table 5.3)}
\label{tab:ea_mapping}
\small
\begin{tabularx}{\textwidth}{*{5}{>{\centering\arraybackslash}X}}
\toprule
\textbf{EA1} & \textbf{EA2} & \textbf{EA3} & \textbf{EA4} & \textbf{EA5} \\
\midrule
Range of Resources & Level of Interaction & Innovation & Consequences & Familiarity \\
\midrule
\checkmark & \checkmark & \checkmark & \checkmark & \checkmark \\
\bottomrule
\end{tabularx}
\end{table}

\begin{itemize}
    \item \textbf{EA1 (Range of Resources)}: Orchestration of diverse software packages, relational databases, cloud VPS servers, and QR optical peripherals.
    \item \textbf{EA2 (Level of Interaction)}: Multi-role operational interfaces catering to students, library desk operators, and chief administrators.
    \item \textbf{EA3 (Innovation)}: Unified synthesis of spatial shelf navigation, digital PDF repository access, and slot-concurrency seat scheduling.
    \item \textbf{EA4 (Consequences for Society and Environment)}: Eradicates paper usage and promotes fair, conflict-free resource sharing.
    \item \textbf{EA5 (Familiarity)}: Engineering a custom, scalable solution under dynamic real-world campus constraints.
\end{itemize}

\section{Summary}
This chapter established strict compliance with international software and security standards, proved significant environmental and societal benefits, verified an 88.6\% cost reduction over proprietary software, and successfully validated all BAETE accreditation dimensions (EP1--EP7, K1--K8, EA1--EA5).
